% ============================================================
% PACHETE NECESARE ÎN PREAMBULUL main.tex
%
% Copiați blocul de mai jos în main.tex, înainte de \begin{document}.
% Dacă un pachet este deja prezent, nu îl duplicați.
% ============================================================

% ── Encoding și limbă (doar pentru pdfLaTeX; XeLaTeX/LuaLaTeX nu au nevoie) ──
\usepackage[utf8]{inputenc}   % caractere românești + $\rightarrow$ din surse
\usepackage[T1]{fontenc}      % fonturi cu diacritice corecte
\usepackage[romanian]{babel}  % silabisire românească

% ── URL-uri (necesar pentru \url{} din secțiunea 7) ──
% Alegeți una dintre variantele de mai jos (hyperref este preferabil pentru teză):
\usepackage{hyperref}         % activează link-uri PDF + oferă \url{}
% SAU, dacă nu vreți link-uri colorate:
% \usepackage[hidelinks]{hyperref}
% SAU, minim (fără link-uri PDF):
% \usepackage{url}

% ── Tabele lungi (necesar pentru tabelul comparativ din secțiunea 7) ──
\usepackage{longtable}

% ── Diagrame TikZ (necesar pentru secțiunea 6) ──
\usepackage{tikz}
\usetikzlibrary{positioning}        % sintaxa  below=0.5cm of <nod>
\usetikzlibrary{shapes.geometric}   % forma    ellipse (diagrama use-case)
\usetikzlibrary{arrows}             % vârful   >=stealth (toate diagramele)

% ── Opțional dar recomandat pentru tabel mai compact ──
% \usepackage{booktabs}   % \toprule, \midrule, \bottomrule (aspect mai curat)
% \usepackage{array}      % tipuri de coloane extinse (m{}, b{})

% ============================================================
% NOTĂ DESPRE ERORI REZOLVATE ÎN FIȘIERELE .tex
%
% Toate caracterele Unicode → (U+2192) din textul vizibil al secțiunilor
% 01, 03, 04, 05 au fost înlocuite cu $\rightarrow$ (LaTeX math mode).
% Caracterele → din comentarii LaTeX (linii ce încep cu %) nu cauzau erori
% și au rămas neschimbate.
%
% Săgeata cu lungime zero din diagrama de secvență (06_diagrame.tex) a
% fost înlocuită cu un nod-notă \node[note] corect.
% ============================================================
